% !TeX root = ../main.tex
% =================================================================================
\chapter{LaTeX Examples}\label{chap:appendix:latex}
% =================================================================================
% TODO: This chapter shows how to use the template. Remove it (in main.tex) before submission.
\TUMremovebeforesubmission{The appendix "LaTeX Examples" of the template}

This chapter shows how to typeset common elements of a thesis with this template.
Look at the source code of this chapter (\texttt{appendix/b-latex-examples.tex}) to see how it is done.

% ---------------------------------------------------------------------------------
\section{Cross References}\label{sec:latex:references}
% ---------------------------------------------------------------------------------

Use \verb|\cref{...}| to reference figures, tables, sections, etc.
It adds the type automatically, e.g., \cref{fig:latex:plot}, \cref{tab:latex:comparison}, \cref{lst:latex:java}, \cref{alg:latex:euclid}, \cref{eq:latex:euler}, and \cref{chap:introduction}.
Use \verb|\Cref{...}| at the beginning of a sentence.
Several references are combined: \cref{fig:latex:plot,fig:latex:sub}.
Use \verb|\refRQ{...}| for research questions defined with \verb|\defRQ{...}{...}|, e.g., \refRQ{resolve}.
Label prefixes such as \texttt{fig:}, \texttt{tab:}, and \texttt{sec:} help you to keep track of your labels.

% ---------------------------------------------------------------------------------
\section{Citations}\label{sec:latex:citations}
% ---------------------------------------------------------------------------------

Cite with \verb|\cite{...}|~\cite{Lamport1994}, cite several works at once~\cite{Basili1994,Runeson2009}, and add page numbers~\cite[p.~42]{Lamport1994}.
Use \verb|\textcite{...}| if the authors are part of the sentence: \textcite{Runeson2009} recommend \ldots.
Manage your references with a reference manager such as JabRef or Zotero.
For web pages, add the \texttt{urldate} field to state when you last visited them~\cite{ISDO}.

% ---------------------------------------------------------------------------------
\section{Acronyms}\label{sec:latex:acronyms}
% ---------------------------------------------------------------------------------

Define acronyms in \texttt{acronyms.tex} and use them with \verb|\gls{...}|.
The first use prints the long form, e.g., \gls{uml}, all further uses only the short form: \gls{uml}.
Use \verb|\glspl{...}| for the plural (\glspl{pdf}), \verb|\Gls{...}| at the beginning of a sentence, and \verb|\glsxtrlong{...}| for the long form only (\glsxtrlong{tum}).
Only acronyms that are used appear in the list of acronyms, together with the page of their first use.
Do not use \verb|\gls{...}| in the abstract, in headings, or in captions, as the first use would then be recorded there.

% ---------------------------------------------------------------------------------
\section{Figures}\label{sec:latex:figures}
% ---------------------------------------------------------------------------------

Put figures into the \texttt{figures} folder and include them with \verb|\includegraphics|.
Prefer vector graphics (PDF, or TikZ as in \cref{fig:concept:overview}) over raster graphics.
Captions of figures are placed below them.

\begin{figure}[htbp]
  \centering
  \begin{subfigure}[b]{0.45\textwidth}
    \centering
    \includegraphics[width=\linewidth]{example-image-a}
    \caption{First sub figure}\label{fig:latex:sub-a}
  \end{subfigure}
  \hfill
  \begin{subfigure}[b]{0.45\textwidth}
    \centering
    \includegraphics[width=\linewidth]{example-image-b}
    \caption{Second sub figure}\label{fig:latex:sub-b}
  \end{subfigure}
  \caption{Two figures side by side using the \texttt{subcaption} package.}
  \label{fig:latex:sub}
\end{figure}

\Cref{fig:latex:sub-a} is a sub figure of \cref{fig:latex:sub}.

\begin{figure}[htbp]
  \centering
  \begin{tikzpicture}
    \begin{axis}[
        width=0.75\textwidth,
        height=5.5cm,
        xlabel={Number of requests},
        ylabel={Latency (\unit{\milli\second})},
        legend pos=north west,
        grid=major,
        grid style={TUMLightGray},
      ]
      \addplot coordinates {(100,12) (200,18) (400,29) (800,52) (1600,101)};
      \addlegendentry{Baseline}
      \addplot coordinates {(100,10) (200,13) (400,17) (800,24) (1600,39)};
      \addlegendentry{Your approach}
    \end{axis}
  \end{tikzpicture}
  \caption{A plot created with \texttt{pgfplots}. The data can also be read from a CSV file.}
  \label{fig:latex:plot}
\end{figure}

% ---------------------------------------------------------------------------------
\section{Tables}\label{sec:latex:tables}
% ---------------------------------------------------------------------------------

Use the \texttt{booktabs} rules and avoid vertical lines.
Captions of tables are placed above them.
Align numbers at the decimal separator with the \texttt{S} column type of \texttt{siunitx} (see \cref{tab:evaluation:results}).

\begin{table}[htbp]
  \centering
  \caption{Comparison of features using \texttt{tabularx}, which fills the given width.}
  \label{tab:latex:comparison}
  \begin{tabularx}{0.8\textwidth}{@{}Xccc@{}}
    \toprule
    Feature                & Tool A & Tool B & Your approach \\
    \midrule
    Detects modifications  & \cmark & \cmark & \cmark        \\
    Resolves conflicts     & \xmark & \cmark & \cmark        \\
    Works offline          & \xmark & \xmark & \cmark        \\
    \bottomrule
  \end{tabularx}
\end{table}

% ---------------------------------------------------------------------------------
\section{Source Code and Algorithms}\label{sec:latex:code}
% ---------------------------------------------------------------------------------

\Cref{lst:latex:java} shows source code typeset with the \texttt{listings} package.
Short code in the text is written with \verb|\lstinline|, e.g., \lstinline[language=Java]|String s = "Hello";|.
To include a complete source file, use \verb|\lstinputlisting|.

\begin{lstlisting}[language=Java, caption={A Java example}, label={lst:latex:java}, float=htbp]
public class HelloWorld {
  // Prints a greeting
  public static void main(String[] args) {
    System.out.println("Hello, TUM!");
  }
}
\end{lstlisting}

Pseudo code is typeset with the \texttt{algpseudocode} package, see \cref{alg:latex:euclid}.

\begin{algorithm}[htbp]
  \caption{Euclid's algorithm}\label{alg:latex:euclid}
  \begin{algorithmic}[1]
    \Procedure{Euclid}{$a, b$}
      \While{$b \neq 0$}
        \State $r \gets a \bmod b$
        \State $a \gets b$
        \State $b \gets r$ \Comment{$b$ decreases in every iteration}
      \EndWhile
      \State \Return $a$
    \EndProcedure
  \end{algorithmic}
\end{algorithm}

% ---------------------------------------------------------------------------------
\section{Mathematics and Units}\label{sec:latex:math}
% ---------------------------------------------------------------------------------

Inline formulas are written as $a^2 + b^2 = c^2$; numbered equations as in \cref{eq:latex:euler}:
\begin{equation}
  e^{i\pi} + 1 = 0 \label{eq:latex:euler}
\end{equation}
Write numbers and units with \texttt{siunitx}: \num{1234567}, \qty{3.5}{\giga\byte}, \qty{20}{\percent}.

% ---------------------------------------------------------------------------------
\section{Quotes, Notes, and TODOs}\label{sec:latex:misc}
% ---------------------------------------------------------------------------------

Use \verb|\enquote{...}| for \enquote{quotation marks} that match the document language.
Footnotes are added with \verb|\footnote{...}|.\footnote{This is a footnote.}
Use \verb|\todo{...}| for notes to yourself.\todo{This is a TODO note.}
\todo[inline]{Inline TODO notes are also possible. Use \texttt{\textbackslash usepackage[disable]\{todonotes\}} in \texttt{settings.tex} to hide all notes.}
\Cref{fig:latex:missing} shows a placeholder for a figure that is not yet finished.

\begin{figure}[htbp]
  \centering
  \missingfigure[figwidth=0.6\textwidth]{A figure that still has to be created.}
  \caption{Placeholder for a missing figure}
  \label{fig:latex:missing}
\end{figure}
